% ---------------------------------------------------------------------------
% Abbildung: FIR-Filter - Ringpuffer (manuell) gegenüber ausgerollter Summe (Faust)
% benötigt: \usepackage{tikz}
%            \usetikzlibrary{arrows.meta,positioning,calc}
% Einbinden mit: % ---------------------------------------------------------------------------
% Abbildung: FIR-Filter - Ringpuffer (manuell) gegenüber ausgerollter Summe (Faust)
% benötigt: \usepackage{tikz}
%            \usetikzlibrary{arrows.meta,positioning,calc}
% Einbinden mit: % ---------------------------------------------------------------------------
% Abbildung: FIR-Filter - Ringpuffer (manuell) gegenüber ausgerollter Summe (Faust)
% benötigt: \usepackage{tikz}
%            \usetikzlibrary{arrows.meta,positioning,calc}
% Einbinden mit: % ---------------------------------------------------------------------------
% Abbildung: FIR-Filter - Ringpuffer (manuell) gegenüber ausgerollter Summe (Faust)
% benötigt: \usepackage{tikz}
%            \usetikzlibrary{arrows.meta,positioning,calc}
% Einbinden mit: \input{figures/abb_fir_ringpuffer}
% ---------------------------------------------------------------------------
\begin{figure}[htbp]
	\centering
	\begin{tikzpicture}[
		>={Stealth[length=2mm]},
		font=\footnotesize,
		cell/.style = {draw, fill=black!4, minimum size=7mm, inner sep=0pt},
		ttl/.style  = {font=\footnotesize\bfseries, align=center},
		note/.style = {font=\scriptsize, align=center}
		]
		
		% ===== links: Ringpuffer der manuellen Implementierung ==================
		\node[ttl] at (0,2.7) {manuelle Implementierung};
		
		\foreach \i in {0,...,7} {
			\node[cell] (c\i) at ({90 - \i*45}:1.75) {\scriptsize \i};
		}
		\foreach \i in {0,...,7} {
			\draw[->, black!45] ({90 - \i*45 - 16}:1.75) arc ({90 - \i*45 - 16}:{90 - \i*45 - 29}:1.75);
		}
		
		\node[note] (pos) at (0,0) {\texttt{pos\_}};
		\draw[->, thick] (pos) -- (c1);
		
		\node[note] at (0,-2.6)
		{Zugriff auf $x[n-k]$ als \texttt{(pos\_ - k) \& (L-1)}\\
			Filterlänge $L$ ist Zweierpotenz, deshalb\\
			Bitmaske statt Division};
		
		% ===== rechts: die vom Faust-Compiler erzeugte Form =====================
		\node[ttl] at (6.4,2.7) {Faust-generierter Code};
		
		\node[draw, rounded corners=2pt, fill=black!4, align=left,
		inner xsep=6pt, inner ysep=5pt] (unroll) at (6.4,0.6) {%
			$y[n] = b_0\,x[n]$\\
			$\qquad\ +\ b_1\,x[n-1]$\\
			$\qquad\ +\ b_2\,x[n-2]$\\
			$\qquad\ \vdots$\\
			$\qquad\ +\ b_{L-1}\,x[n-L+1]$};
		
		\node[note] at (6.4,-2.6)
		{die Summe steht vollständig im Quelltext,\\
			ohne Schleife über die Koeffizienten;\\
			der erzeugte Code wächst mit $L$};
		
	\end{tikzpicture}
	\caption{Zwei Wege zur selben Faltung. Die manuelle Implementierung hält die
		vergangenen Eingangswerte in einem Ringpuffer und läuft in einer Schleife
		über die Koeffizienten. Der Faust-Compiler schreibt dieselbe Summe
		vollständig aus.}
	\label{fig:fir-ringpuffer}
\end{figure}

% ---------------------------------------------------------------------------
\begin{figure}[htbp]
	\centering
	\begin{tikzpicture}[
		>={Stealth[length=2mm]},
		font=\footnotesize,
		cell/.style = {draw, fill=black!4, minimum size=7mm, inner sep=0pt},
		ttl/.style  = {font=\footnotesize\bfseries, align=center},
		note/.style = {font=\scriptsize, align=center}
		]
		
		% ===== links: Ringpuffer der manuellen Implementierung ==================
		\node[ttl] at (0,2.7) {manuelle Implementierung};
		
		\foreach \i in {0,...,7} {
			\node[cell] (c\i) at ({90 - \i*45}:1.75) {\scriptsize \i};
		}
		\foreach \i in {0,...,7} {
			\draw[->, black!45] ({90 - \i*45 - 16}:1.75) arc ({90 - \i*45 - 16}:{90 - \i*45 - 29}:1.75);
		}
		
		\node[note] (pos) at (0,0) {\texttt{pos\_}};
		\draw[->, thick] (pos) -- (c1);
		
		\node[note] at (0,-2.6)
		{Zugriff auf $x[n-k]$ als \texttt{(pos\_ - k) \& (L-1)}\\
			Filterlänge $L$ ist Zweierpotenz, deshalb\\
			Bitmaske statt Division};
		
		% ===== rechts: die vom Faust-Compiler erzeugte Form =====================
		\node[ttl] at (6.4,2.7) {Faust-generierter Code};
		
		\node[draw, rounded corners=2pt, fill=black!4, align=left,
		inner xsep=6pt, inner ysep=5pt] (unroll) at (6.4,0.6) {%
			$y[n] = b_0\,x[n]$\\
			$\qquad\ +\ b_1\,x[n-1]$\\
			$\qquad\ +\ b_2\,x[n-2]$\\
			$\qquad\ \vdots$\\
			$\qquad\ +\ b_{L-1}\,x[n-L+1]$};
		
		\node[note] at (6.4,-2.6)
		{die Summe steht vollständig im Quelltext,\\
			ohne Schleife über die Koeffizienten;\\
			der erzeugte Code wächst mit $L$};
		
	\end{tikzpicture}
	\caption{Zwei Wege zur selben Faltung. Die manuelle Implementierung hält die
		vergangenen Eingangswerte in einem Ringpuffer und läuft in einer Schleife
		über die Koeffizienten. Der Faust-Compiler schreibt dieselbe Summe
		vollständig aus.}
	\label{fig:fir-ringpuffer}
\end{figure}

% ---------------------------------------------------------------------------
\begin{figure}[htbp]
	\centering
	\begin{tikzpicture}[
		>={Stealth[length=2mm]},
		font=\footnotesize,
		cell/.style = {draw, fill=black!4, minimum size=7mm, inner sep=0pt},
		ttl/.style  = {font=\footnotesize\bfseries, align=center},
		note/.style = {font=\scriptsize, align=center}
		]
		
		% ===== links: Ringpuffer der manuellen Implementierung ==================
		\node[ttl] at (0,2.7) {manuelle Implementierung};
		
		\foreach \i in {0,...,7} {
			\node[cell] (c\i) at ({90 - \i*45}:1.75) {\scriptsize \i};
		}
		\foreach \i in {0,...,7} {
			\draw[->, black!45] ({90 - \i*45 - 16}:1.75) arc ({90 - \i*45 - 16}:{90 - \i*45 - 29}:1.75);
		}
		
		\node[note] (pos) at (0,0) {\texttt{pos\_}};
		\draw[->, thick] (pos) -- (c1);
		
		\node[note] at (0,-2.6)
		{Zugriff auf $x[n-k]$ als \texttt{(pos\_ - k) \& (L-1)}\\
			Filterlänge $L$ ist Zweierpotenz, deshalb\\
			Bitmaske statt Division};
		
		% ===== rechts: die vom Faust-Compiler erzeugte Form =====================
		\node[ttl] at (6.4,2.7) {Faust-generierter Code};
		
		\node[draw, rounded corners=2pt, fill=black!4, align=left,
		inner xsep=6pt, inner ysep=5pt] (unroll) at (6.4,0.6) {%
			$y[n] = b_0\,x[n]$\\
			$\qquad\ +\ b_1\,x[n-1]$\\
			$\qquad\ +\ b_2\,x[n-2]$\\
			$\qquad\ \vdots$\\
			$\qquad\ +\ b_{L-1}\,x[n-L+1]$};
		
		\node[note] at (6.4,-2.6)
		{die Summe steht vollständig im Quelltext,\\
			ohne Schleife über die Koeffizienten;\\
			der erzeugte Code wächst mit $L$};
		
	\end{tikzpicture}
	\caption{Zwei Wege zur selben Faltung. Die manuelle Implementierung hält die
		vergangenen Eingangswerte in einem Ringpuffer und läuft in einer Schleife
		über die Koeffizienten. Der Faust-Compiler schreibt dieselbe Summe
		vollständig aus.}
	\label{fig:fir-ringpuffer}
\end{figure}

% ---------------------------------------------------------------------------
\begin{figure}[htbp]
	\centering
	\begin{tikzpicture}[
		>={Stealth[length=2mm]},
		font=\footnotesize,
		cell/.style = {draw, fill=black!4, minimum size=7mm, inner sep=0pt},
		ttl/.style  = {font=\footnotesize\bfseries, align=center},
		note/.style = {font=\scriptsize, align=center}
		]
		
		% ===== links: Ringpuffer der manuellen Implementierung ==================
		\node[ttl] at (0,2.7) {manuelle Implementierung};
		
		\foreach \i in {0,...,7} {
			\node[cell] (c\i) at ({90 - \i*45}:1.75) {\scriptsize \i};
		}
		\foreach \i in {0,...,7} {
			\draw[->, black!45] ({90 - \i*45 - 16}:1.75) arc ({90 - \i*45 - 16}:{90 - \i*45 - 29}:1.75);
		}
		
		\node[note] (pos) at (0,0) {\texttt{pos\_}};
		\draw[->, thick] (pos) -- (c1);
		
		\node[note] at (0,-2.6)
		{Zugriff auf $x[n-k]$ als \texttt{(pos\_ - k) \& (L-1)}\\
			Filterlänge $L$ ist Zweierpotenz, deshalb\\
			Bitmaske statt Division};
		
		% ===== rechts: die vom Faust-Compiler erzeugte Form =====================
		\node[ttl] at (6.4,2.7) {Faust-generierter Code};
		
		\node[draw, rounded corners=2pt, fill=black!4, align=left,
		inner xsep=6pt, inner ysep=5pt] (unroll) at (6.4,0.6) {%
			$y[n] = b_0\,x[n]$\\
			$\qquad\ +\ b_1\,x[n-1]$\\
			$\qquad\ +\ b_2\,x[n-2]$\\
			$\qquad\ \vdots$\\
			$\qquad\ +\ b_{L-1}\,x[n-L+1]$};
		
		\node[note] at (6.4,-2.6)
		{die Summe steht vollständig im Quelltext,\\
			ohne Schleife über die Koeffizienten;\\
			der erzeugte Code wächst mit $L$};
		
	\end{tikzpicture}
	\caption{Zwei Wege zur selben Faltung. Die manuelle Implementierung hält die
		vergangenen Eingangswerte in einem Ringpuffer und läuft in einer Schleife
		über die Koeffizienten. Der Faust-Compiler schreibt dieselbe Summe
		vollständig aus.}
	\label{fig:fir-ringpuffer}
\end{figure}
